\begin{table}[H]
\scriptsize
\setlength{\tabcolsep}{2.5pt}
\renewcommand{\arraystretch}{0.95}

\caption{\label{tab:country-heterogeneity-counterfactual-gap}No-policy counterfactual inflation and inflation inequality (20-country euro area), June 2021--June 2023}
\centering
\resizebox{\textwidth}{!}{%
\begin{tabular}{lrrrrrr}
\toprule
Country & \shortstack{Counterfactual\\ cumulative inflation} & \shortstack{Effect of fiscal policies\\ on average inflation} & \shortstack{Q1 counterfactual\\ cumulative inflation} & \shortstack{Q5 counterfactual\\ cumulative inflation} & \shortstack{Counterfactual\\ inflation gap} & \shortstack{Effect of fiscal policies\\ on inflation inequality}\\
\midrule
Austria & 18.6 & -1.4 & 19.3 & 18.2 & 1.1 & -0.7\\
Cyprus & 12.0 & -0.0 & 12.8 & 12.1 & 0.7 & -0.0\\
Estonia & 33.2 & -0.2 & 39.0 & 31.4 & 7.6 & -0.1\\
Finland & 12.6 & -0.1 & 11.6 & 12.9 & -1.2 & 0.1\\
France & 13.3 & -1.1 & 13.6 & 13.0 & 0.6 & -0.4\\
\addlinespace
Germany & 16.9 & -1.3 & 18.4 & 16.2 & 2.2 & -0.9\\
Greece & 16.6 & -1.9 & 17.8 & 15.8 & 2.0 & -1.4\\
Ireland & 15.2 & -0.3 & 17.6 & 13.6 & 4.1 & -0.2\\
Italy & 16.6 & -0.8 & 18.2 & 15.8 & 2.4 & -0.3\\
Latvia & 29.6 & -0.7 & 37.2 & 27.0 & 10.2 & -0.6\\
\addlinespace
Lithuania & 31.7 & -1.2 & 35.5 & 29.8 & 5.7 & -1.1\\
Luxembourg & 11.9 & -0.5 & 11.6 & 11.5 & 0.1 & -0.3\\
Netherlands & 18.9 & -2.0 & 20.3 & 18.0 & 2.3 & -1.1\\
Portugal & 16.1 & -2.0 & 15.5 & 17.0 & -1.5 & -0.8\\
Slovakia & 30.1 & -4.8 & 33.5 & 27.8 & 5.7 & -3.0\\
\addlinespace
Spain & 13.8 & -2.1 & 14.5 & 13.4 & 1.0 & -0.9\\
\textbf{Euro area} & \textbf{16.0} & \textbf{-1.3} & \textbf{17.0} & \textbf{15.4} & \textbf{1.6} & \textbf{-1.1}\\
\bottomrule
\end{tabular}%
}
\begin{flushleft}
\footnotesize{Notes. Counterfactual cumulative inflation uses the Eurostat-style chained Laspeyres aggregate: COICOP indices are unchained, aggregated with annual HICP weights, and rechained. Inflation gaps are Quintile 1 minus Quintile 5. Counterfactual indices use the same no-policy COICOP component paths reported in the counterfactual appendix; observed 2019 HICP data are used only to chain the January 2020 starting values. They replace observed prices for the affected components with estimated no-policy paths and otherwise retain observed prices; countries without a simulated measure are treated as unchanged. Household-group aggregation uses RAS expenditure weights with the updated HBS--HICP imputed-rent bridge at COICOP level 2 for every country, including Italy and Spain. The euro-area aggregate uses the same annual HICP country weights for all 20 countries as the observed indices. The effect on average inflation is observed cumulative inflation minus counterfactual cumulative inflation. The effect on inflation inequality is the observed gap minus the counterfactual gap. All numeric columns are expressed in percentage points. Sources: Eurostat, national statistical institutes, authors' calculations.}
\end{flushleft}
\end{table}
